\chapter{Conclusion & Future Work}

\section{Limitations and Future Work}
While AniFlowFormer-T advances optical flow estimation for stylized animation, several limitations remain and motivate future directions.

\begin{itemize}
\item \textbf{Very long-range consistency.}
Although the Latent Cost Memory improves temporal stability over short and mid-range horizons, maintaining precise correspondences over hundreds of frames is still challenging. In long shots, small pose variations and repeated patterns can gradually accumulate drift, especially when key structures temporarily disappear due to occlusions or camera motion.

```
\item \textbf{High-frequency stylistic effects.}
Animation frequently employs non-physical effects such as motion streaks, smear frames, impact frames, and intentional shape exaggerations that violate typical correspondence assumptions. These effects can create transient appearances that are inconsistent across time, which may confuse temporal cost construction and lead to unstable matches.

\item \textbf{Under-utilization of region-level supervision.}
AnimeRun provides region matching annotations in addition to dense flow, but the current framework does not explicitly leverage segment- or region-level constraints. Incorporating region correspondences (e.g., character parts or object segments) may further improve contour stability and identity preservation under occlusion.

\item \textbf{Scalability to high-resolution production content.}
Real-world animation pipelines often operate at 2K--4K resolution, where dense temporal matching becomes expensive. While our current implementation remains efficient at 384$\times$768, scaling to higher resolutions may require hierarchical temporal windows, sparse propagation, or adaptive token pruning to reduce compute and memory overhead.
```

\end{itemize}

Future work will explore (i) integrating segment-matching supervision to strengthen boundary and part-level consistency, (ii) extending the temporal horizon via memory compression or retrieval-based mechanisms, and (iii) incorporating animation-specific priors (e.g., smear-frame detection or effect-aware matching) to improve robustness under stylized motion artifacts.

\section{Conclusion}
We presented AniFlowFormer-T, an unsupervised multi-frame temporal Transformer tailored for optical flow estimation in 2D animation. Unlike natural videos, animation exhibits large flat-color regions, sharp contour boundaries, and stylized motion patterns that challenge conventional two-frame models and photometric self-supervision. AniFlowFormer-T addresses these issues by combining a temporal cost volume with a Latent Cost Memory for long-range correspondence aggregation and a Global Temporal Regressor for sequence-level fusion, enabling stable and contour-aware motion estimation without relying on dense ground-truth flow during training.

Quantitative experiments on AnimeRun demonstrate consistent improvements over strong baselines, reducing the overall EPE from 4.21 (RAFT) to \textbf{3.52} and substantially lowering errors under large displacements. On ATD-12K, evaluated through the AnimeInterp frame interpolation protocol, our predicted flows yield higher-quality intermediate frames, improving PSNR from 31.01 dB to \textbf{31.54 dB}. Together, these results underscore the importance of multi-frame temporal reasoning and long-range memory for robust correspondence in stylized animation. We believe AniFlowFormer-T provides a practical foundation for future work in animation analysis and production-oriented automation, including interpolation, editing, and temporally consistent generation.
